磁驱动通过磁场及其空间梯度向柔性体施加远程力和力矩 \citep{abbott,kim}。有效的控制模型需要将这些输入贯穿力学降阶、辨识和数值实现，并在各环节保持同一个动力学对象。本文针对平面夹持--自由磁性导管模型研究这一完整联系。应用背景允许忽略重力；当前基准则限定为小斜率、无接触且不承受周围流体载荷。向更一般情形扩展时，仍需从一般杆力学出发 \citep{simo,antman}。

非线性目标保留完整的耦合余弦磁势。本文不通过消去其高阶项将系统变成独立的线性振子，因为这些项正是持久性分析中非线性频率变化的来源。控制器围绕该标称目标，补偿已辨识的保守残差和比例阻尼。因此，核心问题是：同一学习补偿能否与维度正确的差分格式、物理磁指令，以及对所得离散动力学的严格陈述相衔接。

\subsection{三个相互衔接的子问题}
\label{fw:three-questions}
\textbf{P1：推导并实现差分格式。} 从相容的 Birkhoff 表示出发，哪些端点变量允许定义生成函数，由此得到什么矩阵更新？力学质量矩阵、相流 Jacobian 与端点变换具有不同的维度。将外观类似标量的公式实现为多模态算法时，必须区分这些对象。对于比例阻尼，加权保结构映射又不同于直接作用在有阻尼向量场上的普通中点法。本文明确给出这些构造，并在进入非线性导管模型之前检验非交换矩阵运算。

\textbf{P2：辨识并物理实现残差补偿。} 能否利用状态观测确定保守残差和阻尼估计，同时保留标称的耦合磁势？辨识、控制器和仿真对象必须使用同一坐标，而学习控制器不能读取生成合成数据时所用的未知系数。第二项要求来自物理实现：两个广义模态指令必须先经过统一的力--力矩接口，再重构为局部磁场--梯度。较小的辨识残差本身既不验证激励充分性，也不验证硬件可行性。

\textbf{P3：建立受控格式的结构与收敛。} 将反馈放入同一个中点方程之后，该格式逼近哪个向量场，辨识误差下还保留什么结构，非线性求解又如何影响收敛？有限时间误差必须相对于同一实际闭环计算；实际闭环与标称目标的偏差属于另一项模型误差。自治保守分支可以具有精确辛性，而离散 KAM 还要求解析域、非共振、扭转和定量小量条件。这些推论不能仅凭相图或较小的能量漂移得出。

\subsection{贯穿全文的同一物理与统计模型}
基准由文献材料参考值、磁体制造商规格和明确声明的几何及阻尼选择共同建立。Sylgard 184 参考值对应其发表文献中的配比与固化条件 \citep{cho2021pdms}，不代表当前梁组件的实测值。在求解含磁偏置的特征值问题之前，端部磁芯的平动和转动惯量已纳入模型。这些输入决定三个子问题共同使用的模态尺度、物理解码器和磁输入映射。参数来源见第~\ref{sec:derivation} 节和补充材料~\ref{sec:supp-mechanics}。

力学推导链与学习假设类保持明确区分。哈密顿神经网络学习标量能量，并由其导数决定向量场 \citep{greydanus}；单独表示耗散与驱动，才能保留有意义的功率平衡 \citep{zhong,desai}。复坐标用于组织实位置和动量通道，复特征也可提供具有结构的表示 \citep{trabelsi}。实能量同时依赖复坐标及其共轭；KAM 所需的解析延拓则将底层复化坐标视为独立变量。这一记号或通用逼近定理均不能为任意拟合网络提供经过验证的导数误差界 \citep{hornik,zhang}。

实际实现的实验固定已知动能，通过状态的积分观测联合学习解析势能残差和一个标量阻尼系数。实验使用一个明确声明的磁势特征，从而分别检验辨识误差、数值误差和输入重构。更大的解析读出仍属于理论框架，但本文计算不声称完成不受此限制的隐藏层训练，也不声称给出新的复网络与实网络性能比较。即使在单特征检验中，完整的标称余弦势也始终保留。

\subsection{几何分析与控制目标}
在几何条件和局部可逆性条件成立时，生成函数可构造经典的辛方法 \citep{feng,fengqin,hairer}。本文采用比例阻尼下的构造性 Birkhoff 表示，推导其端点投影与时间中点更新，并在保守极限下恢复普通哈密顿中点映射。四阶中心生成函数保留其完整修正项。本文将这些经典构造应用于所声明的力学与反馈模型，不主张提出新的通用积分方法。

控制器估计阻尼和势能修正，同时保留 $h_0$ 作为标称目标。实际闭环的保守部分为 $h_c=h_0+\delta R$，其中 $\delta R$ 是残差辨识误差。剩余阻尼必须作为非保守项保留。因此，即使精确辛性和经典 KAM 不成立，时间离散格式仍可能收敛。第~\ref{sec:analysis} 节将这一区分表述为受控数值分析结论，并明确包含隐式求解残差的累积效应。

磁势的四次正规形给出可计算的扭转，作为条件性持久性论证的基础。从学习到 KAM 的结论需要超出较小经验损失的假设 \citep{chen2022}，因此本文保留确定性的系数误差及解析误差预算。连续 KAM 为保守目标提供背景 \citep{poschel}，主要的离散结论则使用 \citet{shang} 的生成映射框架。有限振幅的后验认证还需要具体嵌入、不变性缺陷和几何界 \citep{haro2019aposteriori}。本研究中的数值诊断尚未给出这样的新证书。

相关磁机器人研究将降阶模型、视觉反馈和仿真接入物理控制 \citep{francescon2025shape,tong2025navigation}。本文集中研究引入这些附加感知与硬件层之前所需的解析和计算联系。学习一个保守参考并对其进行跟踪也是合理目标，但它产生不同的误差系统。该相关路线保留在补充材料~\ref{sec:supp-tracking}；主文始终围绕残差补偿目标展开。

\subsection{贡献与全文安排}
第一项贡献是从 Birkhoff 一形式、端点投影到矩阵实现的完整推导，并明确物理、模态和相空间的各级维度。第二项贡献是在耦合标称模型上，为残差补偿建立共同的积分辨识与磁输入实现构造。第三项贡献是对受控离散格式进行分析，区分有限时间收敛、精确或加权几何结构，以及条件性的离散 KAM 持久性。数值研究围绕相应接口组织，以分别检验这些联系。

\begin{table*}[t]
\centering
\caption{六节正文中的三个子问题。各行共用同一力学基准和补偿目标。}
\label{tab:roadmap}
\begin{tabularx}{\textwidth}{p{0.15\textwidth}YY}
\toprule
子问题 & 文中的构造 & 验证与适用边界\\
\midrule
P1：离散格式 & 第~\ref{sec:derivation} 节：相容 Birkhoff 形式、端点投影、矩阵中点及加权映射 & 第~\ref{sec:analysis} 节：独立矩阵指数参考、非交换消元及结构诊断\\
P2：学习控制 & 第~\ref{sec:derivation} 节：联合积分辨识与残差补偿；第~\ref{sec:application} 节：物理磁端口 & 第~\ref{sec:analysis} 节：同一对象上的四组控制消融；证据来自无噪声受限特征实验\\
P3：受控分析 & 第~\ref{sec:analysis} 节：隐式格式内的反馈与含求解残差的收敛；第~\ref{sec:kam} 节：离散持久性条件 & 固定模型误差与步长加密诊断；完整解析 KAM 预算及器件标定仍是独立要求\\
\bottomrule
\end{tabularx}
\end{table*}

第~\ref{sec:derivation} 节给出共同力学模型，以及离散格式和学习构造。第~\ref{sec:analysis} 节建立受控格式的有限时间性质，并报告三部分验证序列。第~\ref{sec:kam} 节将保守分支接入离散 KAM。第~\ref{sec:application} 节通过磁端口实现同一指令，并说明尚需满足的全梁和净空要求。第~\ref{sec:conclusion} 节总结对 P1--P3 的回答。补充材料提供详细证明、数值协议及相关的跟踪扩展。
